\FloatBarrier\clearpage\section{Recomendaciones y trabajo futuro}
\label{sec:recomendaciones}

A partir de los resultados verificados y de las limitaciones metodológicas identificadas en la discusión, se formulan recomendaciones operativas para la selección de modelos, el manejo del pipeline experimental y las líneas de continuación científica del proyecto.

\subsection{Recomendaciones de selección de arquitectura por escenario}
\begin{itemize}[leftmargin=1.4em,itemsep=4pt]
\item \textbf{Escenario predeterminado de equilibrio y baja latencia (elegir ResNet-18):} se recomienda \textbf{ResNet-18} cuando el despliegue experimental opere en GPU de gama media o recursos restringidos (4~GiB VRAM, FP32), se requiera baja latencia por imagen en \texttt{batch=1} (mediana de \textbf{7.22~ms}), tiempos cortos de entrenamiento (\textbf{127.5~s/época}) o mayor resiliencia relativa frente a desenfoque (\texttt{blur\_1}, \texttt{blur\_2}), ruido moderado (\texttt{noise\_002}) y recuperación parcial tras prefiltrado espacial (\texttt{+gaussian\_1}).
\item \textbf{Escenario de datos limpios con prioridad en sensibilidad de melanoma (considerar EfficientNet-B0):} puede considerarse \textbf{EfficientNet-B0} únicamente cuando las capturas dermatoscópicas estén limpias y libres de ruido aditivo de alta frecuencia, exista holgura de memoria (805.7~MiB de VRAM) y de latencia (24.60~ms en \texttt{batch=1}), y el criterio prioritario sea maximizar la sensibilidad limpia de melanoma (\textbf{0.6607}) o el desempeño en \texttt{akiec} (Precision = 0.8571, F1 = 0.7143). Su uso no es recomendable en presencia de ruido sin reentrenamiento específico.
\item \textbf{Escenario de restricción extrema de memoria frente a cobertura minoritaria (evitar MobileNetV3-Small):} debe evitarse \textbf{MobileNetV3-Small} siempre que importe la detección de categorías minoritarias o de alto riesgo (\texttt{mel} recall = 0.4167, \texttt{akiec} recall = 0.1633, \texttt{df} recall = 0.2941) o cuando se busque acelerar la inferencia GPU en \texttt{batch=1} (22.40~ms frente a 7.22~ms en ResNet-18). Su ventaja se restringe exclusivamente al menor consumo de VRAM (211.8~MiB) y tamaño de parámetros (1.53~M).
\end{itemize}

\subsection{Recomendaciones sobre preprocesamiento, filtrado y evaluación}
\begin{itemize}[leftmargin=1.4em,itemsep=4pt]
\item \textbf{Advertencia operativa sobre \texttt{+gaussian\_1} y el compromiso Recall/Precision/FP:} el prefiltrado determinista mediante \texttt{+gaussian\_1} (o \texttt{+median\_3}) sobre imágenes afectadas por \texttt{noise\_005} no debe interpretarse como una recuperación equivalente al dato limpio. Aunque en ResNet-18 eleva el \textit{Recall} de melanoma a \textbf{0.6905} (116/168), simultáneamente reduce la \textit{Precision} de melanoma de \textbf{0.5233} a \textbf{0.3648} y eleva los falsos positivos de \textbf{92} (\texttt{clean}) a \textbf{202}; además, en EfficientNet-B0 la cobertura de \texttt{mel} permanece deprimida (0.1250 con \texttt{+gaussian\_1}). Cualquier aplicación de filtrado en inferencia debe acompañarse del reporte conjunto de \textit{Precision}, \textit{Recall} y conteo TP/FP/FN.
\item \textbf{Preservación estricta de \texttt{split.csv} y control antifuga:} en extensiones o reentrenamientos futuros sobre HAM10000, debe conservarse inalterado el manifiesto fijo \texttt{data/processed/split.csv} (\texttt{split\_sha256: 583779fb...}), manteniendo la agrupación por \texttt{lesion\_id} y duplicados binarios, el cálculo exclusivo de estadísticas ($\mu$, $\sigma$) y pesos de clase $w_c$ sobre \texttt{train}, y la restricción de aumento de datos únicamente a entrenamiento, resguardando la partición \texttt{test} tras la confirmación explícita \texttt{-{}-confirmar-test}.
\item \textbf{Medición empírica de latencia en el hardware objetivo:} no debe inferirse la velocidad real de inferencia a partir del conteo de parámetros o FLOPs teóricos. Es indispensable medir la latencia mediana y P95 de \texttt{forward()} en el mismo hardware, precisión y tamaño de lote objetivo (\texttt{batch=1}), aplicando calentamiento previo y sincronización explícita de CUDA, y reportando el rango de dispersión entre corridas.
\end{itemize}

\subsection{Líneas de trabajo futuro}
\begin{enumerate}[leftmargin=1.5em,itemsep=4pt]
\item \textbf{Evaluación con múltiples semillas y ajuste por arquitectura:} repetir los entrenamientos y pruebas con múltiples semillas independientes (p.\ ej., 3--5 semillas) y explorar recetas de regularización diferenciadas por familia arquitectónica para cuantificar la varianza inter-semilla y estrechar la incertidumbre estadística observada entre ResNet-18 y EfficientNet-B0.
\item \textbf{Entrenamiento robusto mediante aumento con ruido y desenfoque:} incorporar perturbaciones gaussianas y desenfoque controlado dentro de la política de aumento en \texttt{train} para evaluar si las redes adquieren invarianza frente a degradaciones de alta frecuencia sin inducir el incremento severo de falsos positivos observado con \texttt{+gaussian\_1} post-hoc.
\item \textbf{Calibración probabilística y optimización de umbrales clínicos:} estudiar métodos de calibración (como escalado de temperatura en \texttt{val}) y ajuste de umbrales de decisión específicos para \texttt{mel}, analizando curvas de costo clínico entre falsos negativos críticos y sobrecarga por falsos positivos.
\item \textbf{Validación externa multicéntrica y auditoría de equidad:} evaluar los modelos entrenados sobre cohortes dermatoscópicas externas de distintos centros y dispositivos, incorporando identificadores de paciente (\texttt{patient\_id}) y metadatos de fototipo cutáneo para verificar la generalización real fuera del dominio HAM10000.
\end{enumerate}
